\documentclass{amsart}
% For dealing with references we use the comment environment
\usepackage{verbatim}
\newenvironment{reference}{\comment}{\endcomment}
%\newenvironment{reference}{}{}
\newenvironment{slogan}{\comment}{\endcomment}
\newenvironment{history}{\comment}{\endcomment}

% For commutative diagrams we use Xy-pic
\usepackage[all]{xy}

% We use 2cell for 2-commutative diagrams.
\xyoption{2cell}
\UseAllTwocells

% We use multicol for the list of chapters between chapters
\usepackage{multicol}

% This is generally recommended for better output
\usepackage{lmodern}
\usepackage[T1]{fontenc}

% For cross-file-references

% Package for hypertext links:
\usepackage{hyperref}

% For any local file, say "hello.tex" you want to link to please

% Theorem environments.
%
\theoremstyle{plain}
\newtheorem{theorem}[subsection]{Theorem}
\newtheorem{proposition}[subsection]{Proposition}
\newtheorem{lemma}[subsection]{Lemma}

\theoremstyle{definition}
\newtheorem{definition}[subsection]{Definition}
\newtheorem{example}[subsection]{Example}
\newtheorem{exercise}[subsection]{Exercise}
\newtheorem{situation}[subsection]{Situation}

\theoremstyle{remark}
\newtheorem{remark}[subsection]{Remark}
\newtheorem{remarks}[subsection]{Remarks}

\numberwithin{equation}{subsection}

% Macros
%
\def\lim{\mathop{\mathrm{lim}}\nolimits}
\def\colim{\mathop{\mathrm{colim}}\nolimits}
\def\Spec{\mathop{\mathrm{Spec}}}
\def\Hom{\mathop{\mathrm{Hom}}\nolimits}
\def\Ext{\mathop{\mathrm{Ext}}\nolimits}
\def\SheafHom{\mathop{\mathcal{H}\!\mathit{om}}\nolimits}
\def\SheafExt{\mathop{\mathcal{E}\!\mathit{xt}}\nolimits}
\def\Sch{\mathit{Sch}}
\def\Mor{\mathop{\mathrm{Mor}}\nolimits}
\def\Ob{\mathop{\mathrm{Ob}}\nolimits}
\def\Sh{\mathop{\mathit{Sh}}\nolimits}
\def\NL{\mathop{N\!L}\nolimits}
\def\CH{\mathop{\mathrm{CH}}\nolimits}
\def\proetale{{pro\text{-}\acute{e}tale}}
\def\etale{{\acute{e}tale}}
\def\QCoh{\mathit{QCoh}}
\def\Ker{\mathop{\mathrm{Ker}}}
\def\Im{\mathop{\mathrm{Im}}}
\def\Coker{\mathop{\mathrm{Coker}}}
\def\Coim{\mathop{\mathrm{Coim}}}

% Boxtimes
%
\DeclareMathSymbol{\boxtimes}{\mathbin}{AMSa}{"02}

%
% Macros for moduli stacks/spaces
%
\def\QCohstack{\mathcal{QC}\!\mathit{oh}}
\def\Cohstack{\mathcal{C}\!\mathit{oh}}
\def\Spacesstack{\mathcal{S}\!\mathit{paces}}
\def\Quotfunctor{\mathrm{Quot}}
\def\Hilbfunctor{\mathrm{Hilb}}
\def\Curvesstack{\mathcal{C}\!\mathit{urves}}
\def\Polarizedstack{\mathcal{P}\!\mathit{olarized}}
\def\Complexesstack{\mathcal{C}\!\mathit{omplexes}}
% \Pic is the operator that assigns to X its picard group, usage \Pic(X)
% \Picardstack_{X/B} denotes the Picard stack of X over B
% \Picardfunctor_{X/B} denotes the Picard functor of X over B
\def\Pic{\mathop{\mathrm{Pic}}\nolimits}
\def\Picardstack{\mathcal{P}\!\mathit{ic}}
\def\Picardfunctor{\mathrm{Pic}}
\def\Deformationcategory{\mathcal{D}\!\mathit{ef}}

\setlength{\emergencystretch}{3em}
\begin{document}
\title[Stacks Project, Tag 09X3]{659 --- Sheaf pullback satisfies qc descent on locally compact Hausdorff spaces}
\maketitle
\noindent Copyright (C) 2005--2025 Johan de Jong. Stacks Project authors. Source: \href{https://stacks.math.columbia.edu/tag/09X3}{Tag 09X3}.

\noindent Editorial difficulty note (not author text). Difficulty H2: The proof reconstructs compatible stalk elements as actual local sections by finite compact-image neighborhoods, then identifies two sheafifications with the original pullback. This non-open-covering descent argument uses the special qc topology in a substantive way..

\noindent Editorial provenance note (not author text). History: excerpt from revision \texttt{a0444\allowbreak 6e57e\allowbreak c1fbc\allowbreak 252a8\allowbreak 71afc\allowbreak ec775\allowbreak 2fb28\allowbreak 07b14}, sites-cohomology.tex, lines 7674--7761. Reference presentation was adapted; original-author context and core support blocks were retained or added. The mathematical wording is unchanged. Licensed under GNU FDL 1.2 or later; no invariant sections or cover texts. See ../licenses/COPYING. Context blocks reproduce author definitions and supporting results; they are not additional counted problems.

\begin{definition}
\label{definition-covering-LC}
Let $\{f_i : X_i \to X\}$ be a family of morphisms with fixed target
in the category $\textit{LC}$. We say this family is a
{\it qc covering}\footnote{This is nonstandard notation.
We chose it to remind the reader of fpqc coverings of schemes.}
if for every $x \in X$ there exist $i_1, \ldots, i_n \in I$ and
quasi-compact subsets $E_j \subset X_{i_j}$ such that
$\bigcup f_{i_j}(E_j)$ is a neighbourhood of $x$.
\end{definition}

\section{Cohomology on Hausdorff and locally quasi-compact spaces}
\label{section-cohomology-LC}

\noindent
We continue our convention to say ``Hausdorff and locally quasi-compact''
instead of saying ``locally compact'' as is often done in the literature.
Let $\textit{LC}$ denote the category whose objects are Hausdorff and
locally quasi-compact topological spaces and whose morphisms are continuous
maps.



\begin{remark}[Set theoretic issues]
\label{remark-set-theoretic-LC}
The category $\textit{LC}$ is a ``big'' category as its objects form
a proper class. Similarly, the coverings form a proper class.
Let us define the {\it size} of a topological space $X$ to be the
cardinality of the set of points of $X$. Choose a function
$Bound$ on cardinals, for example as in
Sets, Equation (\href{https://stacks.math.columbia.edu/tag/046U}{equation bound (Tag 046U)}).
Finally, let $S_0$ be an initial set of objects of $\textit{LC}$,
for example $S_0 = \{(\mathbf{R}, \text{euclidean topology})\}$.
Exactly as in Sets, Lemma \href{https://stacks.math.columbia.edu/tag/000J}{lemma construct category (Tag 000J)}
we can choose a limit ordinal $\alpha$ such that
$\textit{LC}_\alpha = \textit{LC} \cap V_\alpha$
contains $S_0$ and is preserved under all countable limits and
colimits which exist in $\textit{LC}$. Moreover, if $X \in \textit{LC}_\alpha$
and if $Y \in \textit{LC}$ and
$\text{size}(Y) \leq Bound(\text{size}(X))$, then $Y$ is isomorphic
to an object of $\textit{LC}_\alpha$.
Next, we apply Sets, Lemma \href{https://stacks.math.columbia.edu/tag/000X}{lemma coverings site (Tag 000X)}
to choose set $\text{Cov}$ of qc covering on $\textit{LC}_\alpha$
such that every qc covering in $\textit{LC}_\alpha$ is
combinatorially equivalent to a covering this set.
In this way we obtain a site $(\textit{LC}_\alpha, \text{Cov})$
which we will denote $\textit{LC}_{qc}$.
\end{remark}

\begin{lemma}
\label{lemma-describe-pullback-pi}
Let $X$ be an object of $\textit{LC}_{qc}$. Let $\mathcal{F}$ be a
sheaf on $X$. The rule
$$
\textit{LC}_{qc}/X \longrightarrow \textit{Sets},\quad
(f : Y \to X) \longmapsto \Gamma(Y, f^{-1}\mathcal{F})
$$
is a sheaf and a fortiori also a sheaf on $\textit{LC}_{Zar}/X$.
This sheaf is equal to
$\pi_X^{-1}\mathcal{F}$ on $\textit{LC}_{Zar}/X$ and
$\epsilon_X^{-1}\pi_X^{-1}\mathcal{F}$ on $\textit{LC}_{qc}/X$.
\end{lemma}

\begin{proof}
Denote $\mathcal{G}$ the presheaf given by the formula in the lemma.
Of course the pullback $f^{-1}$ in the formula denotes usual
pullback of sheaves on topological spaces. It is immediate
from the definitions that $\mathcal{G}$ is a sheaf for the Zar
topology.

\medskip\noindent
Let $Y \to X$ be a morphism in $\textit{LC}_{qc}$. Let
$\mathcal{V} = \{g_i : Y_i \to Y\}_{i \in I}$ be a qc covering.
To prove $\mathcal{G}$ is a sheaf for the qc topology it
suffices to show that
$\mathcal{G}(Y) \to H^0(\mathcal{V}, \mathcal{G})$
is an isomorphism, see Sites, Section \href{https://stacks.math.columbia.edu/tag/00W1}{section sheafification (Tag 00W1)}.
We first point out that the map is injective as a qc covering
is surjective and we can detect equality of sections at stalks
(use Sheaves, Lemmas \href{https://stacks.math.columbia.edu/tag/0079}{lemma sheaf subset stalks (Tag 0079)} and
\href{https://stacks.math.columbia.edu/tag/008G}{lemma stalk pullback presheaf (Tag 008G)}). Thus
$\mathcal{G}$ is a separated presheaf on $\textit{LC}_{qc}$
hence it suffices to show that any element
$(s_i) \in H^0(\mathcal{V}, \mathcal{G})$
maps to an element in the image of $\mathcal{G}(Y)$
after replacing $\mathcal{V}$ by a refinement
(Sites, Theorem \href{https://stacks.math.columbia.edu/tag/00WB}{theorem plus (Tag 00WB)}).

\medskip\noindent
Identifying sheaves on $Y_{i, Zar}$ and sheaves on $Y_i$ we find that
$\mathcal{G}|_{Y_{i, Zar}}$ is the pullback of $f^{-1}\mathcal{F}$ under
the continuous map $g_i : Y_i \to Y$. Thus we can choose an open covering
$Y_i = \bigcup V_{ij}$ such that for each $j$ there is an open
$W_{ij} \subset Y$ and a section $t_{ij} \in \mathcal{G}(W_{ij})$
such that $V_{ij}$ maps into $W_{ij}$ and such that
$s|_{V_{ij}}$ is the pullback of $t_{ij}$. In other words,
after refining the covering $\{Y_i \to Y\}$ we may assume there
are opens $W_i \subset Y$ such that $Y_i \to Y$ factors through $W_i$
and sections $t_i$ of $\mathcal{G}$ over $W_i$ which restrict
to the given sections $s_i$. Moreover, if $y \in Y$ is in the image
of both $Y_i \to Y$ and $Y_j \to Y$, then the images $t_{i, y}$
and $t_{j, y}$ in the stalk $f^{-1}\mathcal{F}_y$ agree
(because $s_i$ and $s_j$ agree over $Y_i \times_Y Y_j$).
Thus for $y \in Y$ there is a well defined element $t_y$ of
$f^{-1}\mathcal{F}_y$ agreeing with $t_{i, y}$ whenever $y$
is in the image of $Y_i \to Y$.
We will show that the element $(t_y)$ comes from a global section
of $f^{-1}\mathcal{F}$ over $Y$ which will finish the proof of the lemma.

\medskip\noindent
It suffices to show that this is true locally on $Y$, see
Sheaves, Section \href{https://stacks.math.columbia.edu/tag/007X}{section sheafification (Tag 007X)}. Let $y_0 \in Y$.
Pick $i_1, \ldots, i_n \in I$ and
quasi-compact subsets $E_j \subset Y_{i_j}$ such that
$\bigcup g_{i_j}(E_j)$ is a neighbourhood of $y_0$.
Let $V \subset Y$ be an open neighbourhood of $y_0$ contained
in $\bigcup g_{i_j}(E_j)$ and contained in $W_{i_1} \cap \ldots \cap W_{i_n}$.
Since $t_{i_1, y_0} = \ldots = t_{i_n, y_0}$, after shrinking $V$
we may assume the sections $t_{i_j}|_V$, $j = 1, \ldots, n$ of
$f^{-1}\mathcal{F}$ agree. As $V \subset \bigcup g_{i_j}(E_j)$
we see that $(t_y)_{y \in V}$ comes from this section.

\medskip\noindent
We still have to show that $\mathcal{G}$ is equal to
$\epsilon_X^{-1}\pi_X^{-1}\mathcal{F}$ on $\textit{LC}_{qc}$,
resp.\ $\pi_X^{-1}\mathcal{F}$ on $\textit{LC}_{Zar}$.
In both cases the pullback is defined by taking the presheaf
$$
(f : Y \to X)
\longmapsto
\colim_{f(Y) \subset U \subset X} \mathcal{F}(U)
$$
and then sheafifying. Sheafifying in the Zar topology
exactly produces our sheaf $\mathcal{G}$ and the fact
that $\mathcal{G}$ is a qc sheaf, shows that it works as well
in the qc topology.
\end{proof}
\end{document}
